\chapter{Runtime View (Dynamic View)}

\section{Component description}

\subsection{Initialization}
\label{sec:InitDynamic}

\RequirementTag{ADS_SrvcMgr}{027}\\
The initialization of the RTA-SUM SrvcMgr module requires defining a specific function, Runnable_SrvcMgr_Init(), which we invoke from outside of the SUM SrvcMgr.
The user can schedule the initialization function by following one of the mechanisms below:

\begin{itemize}
    \item As BswM Callout \textit{during} the initialization of the ECU;
    \item As Periodic runnable (Task mapping \textbf{is} required)
\end{itemize}

The initialization function \textbf{shall} implement a mechanism that avoids further initialization after its first call.
{\Satisfies{SRS_SrvcMgr 075}}

\RequirementTag{ADS_SrvcMgr}{025}\\
The inizialization function \textbf{shall} also initialize the calibration state of services to ON, by calling the \textit{SD_InitializeServicesCalibrationState} API, then trigger the calibrating off of services through the use of the \textit{SrvcMgr_CalibrationController} client/server interface operation surveyServiceDisableCalibrations.

{\Satisfies{SRS_SrvcMgr 027}}


The initialization routine \textbf{shall} run before the first evaluation of the PermitExecution.

\subsection{Execution Manager}
\label{sec:ExecManagerDynamic}
\RequirementTag{ADS_SrvcMgr}{004}
\\
The execution manager implements the function \code{Runnable_SrvcMgr_ExecutionManager()}. 
This runnable \textbf{shall} detect the information available at the interface \textit{SUSD_PermitExecution} and  produce the shutdown permission value on the interface \textit{SrvcMgr_ShutdownPermission}.
\\
To be scheduled at the sytem level, the function \code{Runnable_SrvcMgr_ExecutionManager()} \textbf{shall} be visible from the outside of the module.
\\
The \code{Runnable_SrvcMgr_ExecutionManager()} \textbf{shall}:
\begin{enumerate}
    \item Check the initialization status \textit{before} reading the \textit{SUSD_PermitExecution} interface on the dedicated port.
    \item Evaluate the execution conditions for the SrvcMgr sub-modules.
	\item Check the status of the request coming from the SPM.
	\item Check the status of the learning session request.
\end{enumerate}

{\Satisfies{SRS_SrvcMgr 060} \Satisfies{SRS_SrvcMgr 061} \Satisfies{SRS_SrvcMgr 062} \Satisfies{SRS_SrvcMgr 063} \Satisfies{SRS_SrvcMgr 055} \Satisfies{SRS_SrvcMgr 056}}

\RequirementTag{ADS_SrvcMgr}{005}

The function \code{Runnable_SrvcMgr_ExecutionManager()} \textbf{should} be time-based scheduled with a periodicity equals to the configuration parameter. 
\ConfigurationParameterExtRef{Sum_SrvcMgr_ExecutionManagerPeriodicity}

\UMLDiagram{generated/diagrams/Sum_SrvcMgr_ExecutionManagement_StateDiag}{SrvcMgr Execution status Management}
{\Satisfies{SRS_SrvcMgr 060} \Satisfies{SRS_SrvcMgr 061} \Satisfies{SRS_SrvcMgr 062} \Satisfies{SRS_SrvcMgr 063} \Satisfies{SRS_SrvcMgr 055} \Satisfies{SRS_SrvcMgr 056}}

\RequirementTag{ADS_SrvcMgr}{026}\\
If the current state of the SUM\_SrvcMgr is in RUN mode, the \textit{Runnable_SrvcMgr_ExecutionManager()}  \textbf{shall} perform certain actions required for DTC  U19A6 57, U19A6 87 and U19A6 54.
These actions are only performed if the battery voltage, received from SUSD via \textit{SUSD_BattVltgSnsrStsAvl} interface, is available and only once. 
In particular, the Execution Manager function \textbf{shall} retrieve the status of the learning table and the incomplete learning reasons to determine whether 
the status of the DTCs U19A6 54 and U19A6 57 should be set or cleared. Additionally, it \textbf{shall} notify SSFM module of the learning table state, by using \textit{SrvcMgr_SSFM_LearningTableStatus} interface, 
to allow SSFM to evaluate conditions for setting DTC U19A6 87.

{\Satisfies{SRS_SrvcMgr 072} \Satisfies{SRS_SrvcMgr 082} \Satisfies{SRS_SrvcMgr 076}}
\subsection{Learning Session Manager}
\subsubsection{DID 0xE0CC}
\label{sec:Ldid0xe0cc}
\RequirementTag{ADS_SrvcMgr}{007}
\\
SUM\_SrvcMgr obtains the value of the learning table state using the SD API Sd_GetLearningTableState to be provided in the DID 0xE0CC which is a 1 byte.
{\Satisfies{SRS_SrvcMgr 045}}

SUM\_SrvcMgr implements the operations ConditionCheckRead and ReadData triggered by Dcm when this DID is requested. Not particular conditions are required to be verified for this DID.

\UMLDiagram{generated/diagrams/Srvcmgr_LearningSession_DID_E0CC}{DID 0xE0CC: Retrieve Learning Table State Workflow}

\subsubsection{DID 0xE0CD}
\label{sec:didE0CD}
\RequirementTag{ADS_SrvcMgr}{016}
\\
SUM\_SrvcMgr retrieves the list of required client services that are expected to be learned but have not been found or discovered. It also maintains a buffer to track the services that are missing from SSFM.\\
{\Satisfies{SRS_SrvcMgr 046} \Satisfies{SRS_SrvcMgr 034}}

SUM\_SrvcMgr implements the operations ConditionCheckRead, ReadDataLength and ReadData triggered by the Dcm when this DID is requested. The conditions verified by the ConditionCheckRead operation are:
\begin{itemize}
\item Learning Table State
\item Maximum number of missing services to be returned into the Dcm buffer.
\end{itemize}

\UMLDiagram{generated/diagrams/Srvcmgr_LearningSession_DID_E0CD}{DID 0xE0CD: Retrieve Incomplete Learning Table - Client section Workflow}

\subsubsection{DID 0xE0CE}
\label{sec:didE0CE}
\RequirementTag{ADS_SrvcMgr}{017}
\\
SUM\_SrvcMgr retrieves required subscriptions that were expected to be sent and acknowledged but that were not acknowledged (server section).\\

SUM\_SrvcMgr implements the ConditionCheckRead, ReadDataLength and ReadData operations triggered by the Dcm when this DID is requested.
In particular, in the operations:
\begin{itemize}
    \item ConditionCheckRead: the conditions of Learning Table State and the maximum number of missing services, to be returned in the Dcm buffer, are checked.
    \item ReadData: the SUM\_SrvcMgr after checking the maximum number of missing event groups using the \code{Sd\_RetrieveIncompleteLearningReasons} API, 
          retrieves the Dtc value from the SPM module, also using the \code{Sd\_GetNodeAddress} API to obtain the node's IP address.
\end{itemize}

\UMLDiagram{generated/diagrams/Srvcmgr_LearningSession_DID_E0CE}{DID 0xE0CE: Retrieve Incomplete Learning Table - Server section Workflow}

{\Satisfies{SRS_SrvcMgr 046} \Satisfies{SRS_SrvcMgr 077}}


\subsubsection{RID 0x0654: Triggering Learning Session}
\label{subsubsec:RID0x0654}
\RequirementTag{ADS_SrvcMgr}{011}\\
SUM\_SrvcMgr implements the Start Routine of RID 0x0654 responsible of triggering the Learning Session.\\

\paragraph{Enable Condition}
RID 0x0654 execution is controlled by four AND conditions:

\begin{enumerate}
    \item The Vehicle Speed Average Driven must be less than or equal to 3 km/h.
    \item The VSAD-related fault-active status from SSFM must be FALSE.
    \item The Propulsion System Active (PropSysActv) signal must be FALSE.
    \item The PSA-related fault-active status from SSFM must be FALSE.
\end{enumerate}

The data sources used by these checks are selected according to configuration:
\begin{itemize}
    \item COM signal reference:
    \begin{itemize}
        \item same signal/service reference: \textit{RP\_Com\_SrvcMgr\_VeActvProp\_YeDrvnVeh}
        \item different references: \textit{RP\_Com\_SrvcMgr\_VePSAR\_b\_PropSysActv} and \textit{RP\_Com\_SrvcMgr\_YeVSPR\_v\_DrvnVehicle}
    \end{itemize}
    \item SSFM single-port support (\texttt{Sum\_Ssfm\_FxI\_FxFA\_SinglePortsSupport}):
    \begin{itemize}
        \item TRUE: use FxFA interfaces (\textit{RP\_SSFM\_SrvcMgr\_FxFA\_VSAD\_PSA} or \textit{RP\_SSFM\_SrvcMgr\_FxFA\_VSAD}/\textit{RP\_SSFM\_SrvcMgr\_FxFA\_PSA})
        \item FALSE: use FsFA interfaces (\textit{RP\_SSFM\_SrvcMgr\_FsFA\_VSAD\_PSA} or \textit{RP\_SSFM\_SrvcMgr\_FsFA\_VSAD}/\textit{RP\_SSFM\_SrvcMgr\_FsFA\_PSA})
    \end{itemize}
\end{itemize}

If any of the above conditions are not met when RID 0x0654 is requested, the routine shall return the negative response code \texttt{DCM\_E\_CONDITIONSNOTCORRECT} and shall not execute the learning session trigger logic.

\paragraph{Implementation}
The implementation consists of a state machine with valid states DCM\_INITIAL and DCM\_PENDING. In DCM\_INITIAL state:

\begin{itemize}
\item SUM\_SrvcMgr shall notify SUM\_SPM of the learning status "PENDING" by using SrvcMgr_SPM_SdLearningStatusIndication
\item the SD API Sd\_ClearLearningTable is called; 
\item services that belong to a Service Group are started using the Sd_ServiceGroupStart() API;
\item the helper Sum\_Srvcmgr\_SetClientSrvcRequestedRemoteSrvcAvailable() is called; it reuses the same activation handling as Runnable\_SrvcMgr\_ActivateServices, ensuring Service Group members use Sd\_ServiceGroupStart/Stop and non-group services use Sd\_Client/ServerServiceSetState plus Sd\_ConsumedEventGroupSetState for requested clients;
\item server services not in a Service Group are set to AVAILABLE using the Sd_ServerServiceSetState() API;
\item client services not in a Service Group are set to REQUESTED using the Sd_ClientServiceSetState() API;
\item all the Consumed Event Groups associated to each Client Service is set to REQUESTER using the Sd_ConsumedEventGroupSetState() API;
\end{itemize}

If Sd_ClearLearningTable() API returns E_NOT_OK, the SPM is notified of the learning status "FAILED" by using SrvcMgr_SPM_SdLearningStatusIndication.

When the state machine transitions to the DCM\_PENDING state, the learning table write status is checked via Sum\_SrvcMgr\_prv\_CheckLearningTableWriteStatus(). Only after the learning session has concluded (ret == E\_OK or E\_NOT\_OK), the following actions are performed:
\begin{itemize}
    \item Call the SPM_SDLearningStatusIndication when the Learning Session completes in order to inform the SPM
    \item Call the Sd_GetLearningTableState and Sd\_RetrieveIncompleteLearningReasons APIs to retrieve the current learning table state and the number of missing event groups.
    \item Populate the LearningTableStatus to be sent to SSFM via the interface \textit{SrvcMgr_SSFM_LearningTableStatus}, allowing SSFM to evaluate conditions for setting DTC U19A6 87.
    \item Set or Clear the  U19A6 54  and U19A6 57 Dtcs, via the correspoding SUM\_SSFM interfaces;
\end{itemize}
While the learning session is still in progress (ret == E\_PENDING), no DTC evaluation or SPM notification is performed.

The complete workflow is depicted in the figure below. \\
{\Satisfies{SRS_SrvcMgr 043} \Satisfies{SRS_SrvcMgr 072} \Satisfies{SRS_SrvcMgr 076} \Satisfies{SRS_SrvcMgr 080} \Satisfies{SRS_SrvcMgr 081} \Satisfies{SRS_SrvcMgr 082} {\Satisfies{SRS_SrvcMgr 083}}

\label{subsubsec:RID0x0654Flow}
\UMLDiagram{generated/diagrams/Srvcmgr_LearningSession_RID_0654}{RID 0x0654: Trigger Learning Session Workflow}

\subsubsection{RID 0x0612}
\label{sec:RID0x0612}
\RequirementTag{ADS_SrvcMgr}{009}
\\
SUM\_SrvcMgr implements the Start Routine of RID 0x0612 responsible of retrieving the content of the Learning Table w.r.t. the Routine Control Option Records received. \\

The implementation will check the consistency of the control option record received and will return the information requested accordingly. The complete workflow is depicted in the figure below.

\UMLDiagram{generated/diagrams/Srvcmgr_LearningSession_RID_0612}{RID 0x0612: Retrieve Learning Table Workflow}

{\Satisfies{SRS_SrvcMgr 045}}

\paragraph{Enable Conditions}
\RequirementTag{ADS_SrvcMgr}{032}

RID 0x0612 execution is controlled by two conditions that must be satisfied. These signals can be accessed through two different port interfaces depending on the system configuration:

\begin{enumerate}
    \item The Propulsion System Active signal must be FALSE. This signal is read from the required port \textit{RP\_Com\_SrvcMgr\_VePSAR\_b\_PropSysActv} via RTE interface at runtime.
    \item The Vehicle Speed Average Driven must be less than or equal to 3 km/h. This signal is read from the required port \textit{RP\_Com\_SrvcMgr\_YeVSPR\_v\_DrvnVehicle} via RTE interface at runtime.
\end{enumerate}

In case of reference to the same service-based or signal-based (group signal) signals, the SrvcMgr module can alternatively use the port interface \textit{RP\_Com\_SrvcMgr\_VeActvProp\_YeDrvnVeh}. This interface is introduced when both the Propulsion System Active signal and Vehicle Speed Average Driven are provided as service-based signals, allowing unified access to both signals through a single interface for RID 0x0612 and RID 0x0654 enable condition checking.

If either condition is not met when RID 0x0612 is requested, the routine shall return the negative response code \texttt{DCM\_E\_CONDITIONSNOTCORRECT} and shall not execute the learning table retrieval logic.

{\Satisfies{SRS_SrvcMgr 078} \Satisfies{SRS_SrvcMgr 079} \Satisfies{SRS_SrvcMgr 080} \Satisfies{SRS_SrvcMgr 081}}

\subsection{Service Manager}
\label{sec:SrvcManagerDynamic}
\subsubsection{API Runnable_SrvcMgr_ActivateServices()}
\RequirementTag{ADS_SrvcMgr}{022}
\\
This API will be called through the C/S interface SrvcMgr\_SPM\_ActivateServices by the SPM module.
The purpose of this API is to set a list of Service ID (received as input parameter) the state (in function of a Command List received as input parameter).
\begin{itemize}
    \item For services that belong to a Service Group, the SD APIs Sd\_ServiceGroupStart()/Sd\_ServiceGroupStop() are used. In this case SrvcMgr does not call Sd\_ClientServiceSetState or Sd\_ServerServiceSetState for those services.
    \item For services that do not belong to a Service Group, the helper Sum\_SrvcMgr\_HandlerServiceCommand is used:
    \begin{itemize}
        \item It determines if the service is Client or Server and calls the relevant SD API (Sd\_ClientServiceSetState or Sd\_ServerServiceSetState) based on the activationCommandList value (TRUE = activation, FALSE = disabling).
        \item If the SD API returns E\_OK, the commandStatusList entry is set to E\_OK\_LOCAL\_ENABLED (Server Service) or E\_OK\_REMOTE\_ENABLED (Client Service).
        \item For a requested and enabled Client Service that has SdConsumedEventGroup(s), Sd\_ConsumedEventGroupSetState() is called for each associated SdConsumedEventConfiguredEvent. If any Sd\_ConsumedEventGroupSetState() call returns E\_NOT\_OK, the commandStatusList entry is set to E\_GENERAL\_ERROR; otherwise it remains E\_OK\_REMOTE\_ENABLED.
        \item If the SD API returns E\_NOT\_OK, the commandStatusList entry is set to E\_GENERAL\_ERROR.
    \end{itemize}
\end{itemize}
The result of these operations are retrieved by populating an input array Command Status List.
\UMLDiagram{generated/diagrams/SrvcMgr_ServiceManager_ActiveServices}{Runnable_SrvcMgr_ActivateServices workflow.}
{\Satisfies{SRS_SrvcMgr 047} \Satisfies{SRS_SrvcMgr 048} \Satisfies{SRS_SrvcMgr 049} \Satisfies{SRS_SrvcMgr 050} \Satisfies{SRS_SrvcMgr 071} \Satisfies{SRS_SrvcMgr 055} \Satisfies{SRS_SrvcMgr 056}}
{\Satisfies{SRS_SrvcMgr 073}}
\RequirementTag{ADS_SrvcMgr}{024}
\\
Moreover, this API SHALL track if a request from SPM is ongoing in order to prevent that Execution Manager votes for shutdown.
{\Satisfies{SRS_SrvcMgr 055} \Satisfies{SRS_SrvcMgr 056}}
\subsection{Calibration Controller}
\label{subsec:CalibControllerDynamic}

\subsubsection{Runnable\_SrvcMgr\_DisableService}
\RequirementTag{ADS_SrvcMgr}{020}

SUM\_SrvcMgr implements the Runnable\_SrvcMgr\_DisableService that will be called through the C/S interface SrvcMgr\_ServiceCalibrationOff (operation SrvcMgr\_DisableService) 
by the SWCs to change the state (disable) of services through SD APIs:

\begin{itemize}
\item Call to Sd\_ServerServiceSetState()/ClientServiceSetState() (depending on whether the service is client or server);  
\item Call to Sd\_CalibrateServicesOff() (SD extension);
\end{itemize}

The return value of the implementation is based on:

\begin{itemize}
\item The consistency of the Service ID received with respect to the Sd configuration;
\item The return value of the SD APIs;
\end{itemize}

\UMLDiagram{generated/diagrams/SrvcMgr_FlowDiagram_DisableService}{API SrvcMgr\_DisableService workflow} 

{\Satisfies{SRS_SrvcMgr 021} \Satisfies{SRS_SrvcMgr 023} \Satisfies{SRS_SrvcMgr 024}}

\subsubsection{Runnable\_SrvcMgr\_DisableServiceGroup}
\RequirementTag{ADS_SrvcMgr}{021}

SUM\_SrvcMgr implements the Runnable\_SrvcMgr\_DisableServiceGroup that will be called through the C/S interface SrvcMgr\_ServiceCalibrationOff (operation SrvcMgr\_DisableServiceGroup) 
by the SWCs to change the state (disable) of services group through SD APIs:

\begin{itemize}
\item Call to Sd\_ServiceGroupStop;  
\item Call to Sd\_CalibrateServicesOff() (SD extension);
\end{itemize}

The return value of the implementation is based on:

\begin{itemize}
\item The consistency over the Service Group ID received w.r.t. the Sd configuration;
\item The return value of the SD APIs;
\end{itemize}

\UMLDiagram{generated/diagrams/SrvcMgr_FlowDiagram_DisableServiceGroup}{API SrvcMgr\_DisableService workflow}
 
{\Satisfies{SRS_SrvcMgr 022} \Satisfies{SRS_SrvcMgr 025} \Satisfies{SRS_SrvcMgr 026}}

\newpage
\section{Run-time scenarios}
\label{sec:RunTimeScenarios}

\subsection{Initialization}
\RequirementTag{ADS_SrvcMgr}{028}
\UMLDiagram{generated/diagrams/Sum_SrvcMgr_Initialization}{SrvcMgr Initialization}
{\Satisfies{SRS_SrvcMgr 075}}

\subsection{SrvcMgr Execution status Management}
\RequirementTag{ADS_SrvcMgr}{006}
\UMLDiagram{generated/diagrams/Sum_SrvcMgr_ExecutionManager_SeqDiag}{SrvcMgr Execution Manager - Functionality enabler}
{\Satisfies{SRS_SrvcMgr 060} \Satisfies{SRS_SrvcMgr 061} \Satisfies{SRS_SrvcMgr 062} \Satisfies{SRS_SrvcMgr 063}}

\subsection{SrvcMgr Learning Session Management}
\RequirementTag{ADS_SrvcMgr}{018}
\label{sec:RID0x0654LearningPhase}
\UMLDiagram{generated/diagrams/SrvcMgr_SeqDiagram_LearningPhase}{SrvcMgr Learning Session - Learning Phase}
{\Satisfies{SRS_SrvcMgr 043} \Satisfies{SRS_SrvcMgr 072}}

\RequirementTag{ADS_SrvcMgr}{019}
\UMLDiagram{generated/diagrams/SrvcMgr_LrnManager_SeqDiagram_Runtime}{SrvcMgr Learning Session - Runtime}
{\Satisfies{SRS_SrvcMgr 045} \Satisfies{SRS_SrvcMgr 046} \Satisfies{SRS_SrvcMgr 077}  \Satisfies{SRS_SrvcMgr 034}}

\subsection{SrvcMgr Service Manager Management}
\RequirementTag{ADS_SrvcMgr}{023}
\UMLDiagram{generated/diagrams/SUM_SrvcMgr_ServiceManagement}{SrvcMgr Service Manager - ActivateServices}
{{\Satisfies{SRS_SrvcMgr 047} \Satisfies{SRS_SrvcMgr 048} \Satisfies{SRS_SrvcMgr 049} \Satisfies{SRS_SrvcMgr 050} \Satisfies{SRS_SrvcMgr 055} \Satisfies{SRS_SrvcMgr 056} \Satisfies{SRS_SrvcMgr 071}}
}

\section{DTC Management}
\label{sec:DTCs}
\RequirementTag{ADS_SrvcMgr}{030}\\
\textbf{DTC U19A6 57 - Invalid/Incompatible SW Component}\\
The Service Manager shall set the DTC U19A6 57 via \textit{SrvcMgr\_SSFM\_SrvcDscvryFlr\_U19A6\_57} if \textit{Sd\_GetLearningTableState} returns \texttt{LEARNING\_INCOMPLETE} and \\\textit{Sd\_RetrieveIncompleteLearningReasons.num\_Of\_Missing\_EventGroups} $>$ 0.\\
The DTC shall be cleared if \textit{Sd\_GetLearningTableState} returns \texttt{LEARNING\_COMPLETE} or \texttt{LEARNING\_INCOMPLETE} with zero missing event groups.\\
This check is performed in the the Execution Manager, after RID 0x0654 execution, and after Service 0x14 (Clear Diagnostic Information) confirmation via \textit{SSFM\_ClearDiagReq}.

\UMLDiagram{generated/diagrams/SUM_SrvcMgr_SSFM_ClearDiagReq_U19A6_57}{The Service Manager (SrvcMgr) handles DTC U19A6 57 upon receipt the confirmation of Service 0x14 (Clear Diagnostic Information) from SSFM.}

\textbf{DTC U19A6 54 - Missing Setup/Learn} \\
The Service Manager shall set the DTC U19A6 54 via \textit{SrvcMgr\_ERRH\_U19A6\_54\_Status} if \textit{Sd\_GetLearningTableState}
returns \texttt{CLEARED}, otherwise reset the DTC. This action is performed in the the Execution Manager, after RID 0x0654 execution,
and after Service 0x14 (Clear Diagnostic Information) confirmation via \textit{SSFM\_ClearDiagReq}.\\
Note: The interaction with the SSFM module is similar to the diagram shown in the previously DTC.

\textbf{DTC U19A6 87 - Missing Message}\\
If the \textit{Sd\_GetLearningTableState} returns \texttt{LEARNING\_COMPLETE} or \texttt{LEARNING\_INCOMPLETE}, the Service Manager shall forward information on the learning table status to the SSFM module, via the \textit{SrvcMgr\_SSFM\_LearningTableStatus} interface.
This information allows the SSFM to evaluate the conditions for setting DTC U19A6 87. 
This action is performed in the Execution Manager and after RID 0x0654 execution.

{\Satisfies{SRS_SrvcMgr 076} \Satisfies{SRS_SrvcMgr 072} \Satisfies{SRS_SrvcMgr 082} } 

